% Public manuscript source. Build with XeLaTeX and Biber.
\documentclass[11pt]{amsart}

\usepackage{amsmath, amsthm, amssymb, amsfonts, mathtools, mathrsfs}
\usepackage{hyperref}
\usepackage[backend = biber, style = alphabetic]{biblatex}

\addbibresource{references.bib}

\theoremstyle{plain}
\newtheorem{theorem}{Theorem}[section]
\newtheorem{lemma}[theorem]{Lemma}
\newtheorem{corollary}[theorem]{Corollary}
\newtheorem{proposition}[theorem]{Proposition}

\theoremstyle{definition}
\newtheorem{example}[theorem]{Example}

\theoremstyle{remark}
\newtheorem{remark}[theorem]{Remark}

\renewcommand{\proofname}{Proof}
\setlength{\emergencystretch}{2em}

% Equation numbers restart in each section.
\numberwithin{equation}{section}

\title[Cohomological Milnor Fibrations]{Cohomological Milnor Fibrations of Landau--Ginzburg Models}

\author{Huijun Fan}
\address{School of Mathematics and Statistics, Wuhan University, Wuhan, Hubei, China}
\email{fanhuijun@whu.edu.cn, fanhj@math.pku.edu.cn}
\author{Jieran Wang}
\address{School of Mathematics, Sun Yat-sen University, Guangzhou, Guangdong, China}
\email{wangjr88@mail.sysu.edu.cn}
\date{Working draft, version~0.1.0 (October 17, 2026).}

\keywords{Landau--Ginzburg model; Milnor fibration; Milnor number; tame polynomial; Gauss--Manin connection; $tt^*$-geometry}
\subjclass[2020]{32S25; 32S40; 14N35; 32Q15}

\begin{document}

\begin{abstract}
    We study the global topology of \emph{fundamentally tame} Landau--Ginzburg (LG) models on complete K\"ahler manifolds.
    We first show that the superpotential $f$ is surjective onto $\mathbb{C}$: in the non-proper holomorphic setting this is far from automatic, and here it is deduced from fundamental tameness together with the completeness of $M$.
    We then prove that $f$ restricts to a locally trivial fibration over the set of its regular values, which yields a global generalization of the classical Milnor fibration and provides an Ehresmann-type fibration in the non-proper setting.
    Assuming that $f$ has finitely many critical points, we use a Morse-theoretic argument to show that the relative homology of $\bigl(M,f^{-1}(c)\bigr)$, for a regular value $c$, is concentrated in the middle degree and has rank equal to the global Milnor number $\mu_f$.
    We further identify the relative homotopy type of the fibers with that of the sublevel sets of $f$ at infinity.
    Finally, for deformations satisfying a uniform tameness condition on a Stein manifold, we construct the relative cohomology bundles at positive and negative infinity, each equipped with a natural flat Gauss--Manin connection.
    Their direct sum is called the \emph{cohomological Milnor bundle}.
\end{abstract}

\maketitle

\section{Introduction}

Landau--Ginzburg (LG) B-models and $tt^*$-geometry occupy a central place in mathematical physics and singularity theory, and their local theory is by now classical.
Locally one studies a holomorphic function with an isolated singularity at the origin: the basic invariants of the model are supplied by the Milnor fibration together with the Milnor number $\mu_f$ \cite{milnor1968singular}, by the theory of primitive forms and Frobenius manifolds \cite{Saito1981Primitive,SaitoTakahashi2008,Hertling2003}, and by the Hodge-theoretic description through the Brieskorn lattice \cite{saito1989brieskorn}, while the topological-field-theoretic background is provided by \cite{CecottiVafa1991,Dubrovin1993,Dubrovin1996}.
The analytic and deformation-theoretic framework underlying the present work was developed by Fan \cite{fan2011,FanFang2021}, and the $L^2$-Hodge theory of the twisted Cauchy--Riemann operators of LG models was established by Li--Wen \cite{LI2022108165} and Wen \cite{Wen2016,WenFan2015}; see also \cite{FanJarvisRuan2013}.

The passage from the local picture to the global one is, however, not formal, and two difficulties arise.
First, it is not evident whether $f$ is surjective onto $\mathbb{C}$.
Second, the classical proof of the Milnor fibration over a disk invokes the Ehresmann fibration theorem, whose argument requires the map to be proper; without properness the flow that lifts a vector field from the base need not be complete, and the trivialization argument breaks down.
For polynomial functions the corresponding global theory was developed by Broughton \cite{Broughton1988}, where tameness of the polynomial controls the behavior at infinity; see also \cite{arnold2012singularities,Batyrev1993,dimca1992singularities,Kulikov_1998,LeDungTrang1977} for the singularity-theoretic background and \cite{Dubrovin1990ModernIII} for the homology theory involved.

The purpose of this article is to establish the analogous global structure for \emph{fundamentally tame} LG models on complete K\"ahler manifolds, within the analytic framework of \cite{fan2011,LI2022108165,Wen2016}.
Throughout, $M$ is a connected, noncompact, complete K\"ahler manifold of complex dimension $n$, and $f:M\to\mathbb{C}$ is a nonconstant holomorphic function.
It is fundamental tameness, together with the completeness of $M$, that replaces properness in the two obstructions above.
In the form in which it is used below, fundamental tameness provides a compact subset $K\subset M$ and a constant $C>0$ such that
\begin{equation*}
    \|\nabla f(p)\| > C \qquad\text{for all } p\in M\setminus K,
\end{equation*}
so that all critical points of $f$ lie in the compact set $K$, while the behavior of $f$ away from $K$ is controlled by a uniform gradient bound.

Our first result is a \emph{range theorem} (Theorem~\ref{thm:range}): if $(M,f)$ is a fundamentally tame LG model, then $f(M)=\mathbb{C}$.
Its proof is elementary, yet it already exhibits the mechanism that will be used throughout.
Writing $u=\operatorname{Re} f$ and $v=\operatorname{Im} f$, one shows that $u$ is unbounded both above and below by following the negative gradient flow of $u$: as the flow cannot remain in any compact set, completeness of $M$ forces its length to tend to infinity, while the gradient bound forces $u$ to decrease at a uniform rate.
The same argument applied to the level sets of $u$ then shows that $v$ is unbounded along $u^{-1}(\alpha)$ for every $\alpha\in\mathbb{R}$, which yields the surjectivity of $f$.

Our second result is a \emph{fibration theorem} (Proposition~\ref{prop:fibration}): $f$ restricts to a locally trivial fibration $f^{-1}\bigl(f(M)\setminus D\bigr)\to f(M)\setminus D$ over the set of its regular values, where $D$ denotes the set of critical values of $f$.
The argument follows the strategy of the Ehresmann fibration theorem and is inspired by \cite[Lemma 4.3]{Broughton1988}; since $f$ is not assumed to be proper, the lifted vector field is constructed explicitly, and fundamental tameness together with the completeness of $M$ is used to guarantee global existence of its flow.
This yields a global generalization of the classical Milnor fibration and provides an Ehresmann-type fibration in the non-proper setting.

Our third result computes the relative homology of the smooth fibers (Theorem~\ref{thm:main}): if $(M,f)$ is a fundamentally tame LG model whose superpotential has finitely many critical points, then for every regular value $c$ of $f$ the relative homology groups $H_*\bigl(M,f^{-1}(c)\bigr)$ are concentrated in degree $n$, with
\begin{equation*}
    H_n\bigl(M,f^{-1}(c)\bigr)\;\cong\;\mathbb{Z}^{\mu_f},
\end{equation*}
where $\mu_f=\sum_{p\in\operatorname{Crit}(f)}\mu_f(p)$ is the global Milnor number of $f$.
The proof adapts Milnor's classical Morse-theoretic argument \cite{milnor1963morse} to the non-proper, noncompact situation at hand: the auxiliary function $F_a=|f_a-c|^2$ is a nonnegative function whose zero locus is $f^{-1}(c)$ and which, outside a neighborhood of this zero locus, has exactly $\mu_f$ nondegenerate critical points, all of Morse index $n$, while compactness of the sublevel sets is replaced by the completeness of $M$ together with tameness of $F_a$ away from that neighborhood, in the spirit of Broughton \cite{Broughton1988}.
We also clarify the relation between regular fibers and the sublevel sets of $f$ at infinity (Proposition~\ref{prop:homotopy}): for $|\operatorname{Re} c|$ large, the pair $\bigl(M,f^{-1}(c)\bigr)$ is homotopy equivalent to the pair formed by $M$ and the corresponding sublevel set, and the common relative homotopy type is independent of $c$.

Our fourth and last result concerns deformations of the superpotential.
For a deformation $F(z,t)=f(z)+\sum_{i=1}^{m}t^i g_i(z)$ over a bounded parameter domain $S\subset \mathbb{C}^m$, and under a uniform tameness assumption (Theorem~\ref{thm:deformation}), the level set $E_c=\{F=c\}$ is a smooth hypersurface of $M\times S$ for $|\operatorname{Re} c|$ sufficiently large, and the projection $\pi:E_c\to S$ is a locally trivial smooth fibration.
If $M$ is Stein, combining this with the homotopy type at infinity yields (Theorem~\ref{thm:bundle}) the positive and negative relative cohomology bundles $\mathcal{H}^{+}$ and $\mathcal{H}^{-}$, each of rank $\mu$ and equipped with a natural flat Gauss--Manin connection.
We call their direct sum $\mathcal{H}_{\mathrm{tot}}=\mathcal{H}^{+}\oplus\mathcal{H}^{-}$ the \emph{cohomological Milnor bundle} associated with the deformation $F$.
In contrast with the classical Picard--Lefschetz picture, in which the Lefschetz thimbles of a Morse function form a basis of the relative homology, the construction above does not require $f_t$ to be Morse, and both bundles and their flat connections are defined over the whole deformation parameter space.

The article is organized as follows.
Section~\ref{sec:range} proves the range theorem, together with the auxiliary lemma on maximal integral curves of $-\nabla u$ that the argument requires, and records an application to holomorphic embeddings and polynomial superpotentials.
Section~\ref{sec:fibration} establishes the fibration theorem.
Section~\ref{sec:relative} computes the relative homology of the smooth fibers, and Section~\ref{sec:construction} constructs the cohomological Milnor bundles and relates the regular fibers to the sublevel sets of $f$ at infinity.

\section{Range Theorem}\label{sec:range}

As a preparation for our construction, we establish some basic properties of fundamentally tame LG models, beginning with the range theorem for $f$.
This result will later allow us to identify the relative homotopy type of the sublevel sets of $f$ with that of the fibers of $f$ when the real part of $f$ is sufficiently large.

\begin{theorem}\label{thm:range}
    If $(M, f)$ is a fundamentally tame LG model, then $f(M) = \mathbb{C}$.
\end{theorem}

\begin{proof}
    Let $u = \operatorname{Re} f$ and $v = \operatorname{Im} f$.
    Then $u$ and $v$ are nonconstant harmonic functions on $M$.

    \medskip
    \noindent\emph{Step 1: $u(M) = \mathbb{R}$.}

    Since $M$ is connected, the image $u(M)$ is an interval.
    To prove $u(M) = \mathbb{R}$, it therefore suffices to show that
    \begin{equation}
        \inf_{p \in M} u(p) = -\infty \quad \text{and} \quad \sup_{p \in M} u(p) = +\infty.
    \end{equation}
    We prove the statement for the infimum; the argument for the supremum is analogous.

    Assume for contradiction that
    \begin{equation}
        \alpha_m := \inf_{p \in M} u(p) > -\infty.
    \end{equation}

    Since $f$ is fundamentally tame, there exist a compact subset $K \subset M$ and a positive constant $C$ such that
    \begin{equation}
        \|\nabla f(p)\| > C \quad \text{for all } p \in M \setminus K.
    \end{equation}
    In particular, the critical locus of $f$ is contained in $K$.

    Let
    \begin{equation}
        \beta_m := \min_{p \in K} u(p).
    \end{equation}

    By the maximum principle, $u$ does not attain its minimum on $M$.
    Hence
    \begin{equation}
        \alpha_m < \beta_m,
    \end{equation}
    and consequently
    \begin{equation}
        (\alpha_m, \beta_m) \subset u(M).
    \end{equation}
    Moreover,
    \begin{equation}
        u^{-1}(\alpha_m, \beta_m) \subset M \setminus K.
    \end{equation}

    Choose $p \in M$ such that $u(p) \in (\alpha_m, \beta_m)$ and consider the negative gradient flow
    \begin{equation}
        \begin{cases}
            \phi'(t) = -\nabla u(\phi(t)),\\
            \phi(0) = p.
        \end{cases}
    \end{equation}

    Let $d$ denote the distance function on $M$.
    For $t \geq 0$ we have $u(\phi(t)) \leq u(p) < \beta_m$, hence $\phi(t) \in M \setminus K$.
    Since $\|\nabla u\| = \frac{1}{2}\|\nabla f\|$, it follows that
    \begin{equation}
        \|\nabla u(\phi(t))\| > \tfrac{1}{2} C.
    \end{equation}

    Since
    \begin{equation}
        u(p) - u(\phi(t)) = \int_0^t -\frac{\mathrm{d}}{\mathrm{d}\tau} u(\phi(\tau)) \, \mathrm{d}\tau = \int_0^t \|\nabla u(\phi(\tau))\|^2 \, \mathrm{d}\tau,
    \end{equation}
    we can estimate
    \begin{equation}\label{eqn:2.11}
        u(p) - u(\phi(t)) \geq \frac{1}{2} C \int_0^t \|\phi'(\tau)\| \, \mathrm{d}\tau \geq \frac{1}{2} C \, d(\phi(t), p),
    \end{equation}
    while on the other hand
    \begin{equation}\label{eqn:2.12}
        u(p) - u(\phi(t)) \geq \frac{1}{4} C^2 t.
    \end{equation}

    Let $[0, T)$ be the maximal forward interval of existence of $\phi(t)$.

    If $T = +\infty$, then by \eqref{eqn:2.12},
    \begin{equation}
        u(p) - u(\phi(t)) \to +\infty \quad \text{as } t \to +\infty,
    \end{equation}
    which contradicts the assumption that $\alpha_m > -\infty$.

    If $T < +\infty$, then by maximality of the interval the curve $\phi(t)$ cannot remain in any compact subset of $M$.
    Since $M$ is complete, this implies
    \begin{equation}
        d(\phi(t), p) \to \infty \quad \text{as } t \to T^-.
    \end{equation}
    By \eqref{eqn:2.11} we again obtain
    \begin{equation}
        u(p) - u(\phi(t)) \to +\infty \quad \text{as } t \to T^-,
    \end{equation}
    which is a contradiction.

    Therefore $\inf_{p \in M} u(p) = -\infty$.
    Applying the same argument to $-u$ gives $\sup_{p \in M} u(p) = +\infty$, and hence
    \begin{equation}
        u(M) = \mathbb{R}.
    \end{equation}

    The argument above in fact yields the following lemma.

    \begin{lemma}\label{lem:integral-curve}
        Let $M$ be a complete manifold.
        If $\gamma : [0, T) \to M$ is a maximal forward integral curve of $-\nabla u$ and $\|\nabla u\| > C$ on $\operatorname{Im} \gamma$ for some $C > 0$, then
        \begin{equation}
            u(\operatorname{Im} \gamma) = (-\infty, u(\gamma(0))].
        \end{equation}
    \end{lemma}

    \medskip
    \noindent\emph{Step 2: $v\big(u^{-1}(\alpha)\big) = \mathbb{R}$ for every $\alpha \in \mathbb{R}$, that is,}
    \begin{equation}
        v\big(u^{-1}(\alpha)\big) = \mathbb{R}.
    \end{equation}

    Since Step 1 shows $u(M) = \mathbb{R}$, the level set $u^{-1}(\alpha)$ is nonempty.

    Because $f$ is holomorphic, the Cauchy--Riemann equations imply
    \begin{equation}
        \nabla v = J(\nabla u), \qquad \langle \nabla u, \nabla v \rangle = 0,
    \end{equation}
    where $J$ denotes the complex structure.
    Hence along any integral curve $\gamma(t)$ of $\pm \nabla v$,
    \begin{equation}
        \frac{\mathrm{d}}{\mathrm{d}t} u(\gamma(t)) = \langle \nabla u, \gamma'(t) \rangle = \pm \langle \nabla u, \nabla v \rangle = 0.
    \end{equation}
    Therefore every integral curve of $\pm \nabla v$ that starts in $u^{-1}(\alpha)$ remains inside this level set.

    Since the critical locus of $f$ is contained in the compact set $K$, it is a compact complex analytic space.
    Consider its connected components.
    By the maximum principle for holomorphic functions on connected complex analytic spaces \cite[Chapter 5, Section 5]{GrauertRemmert2004}, the function $f$ is constant on each connected component, so that on each such component both $u$ and $v$ are constant.
    Each component is therefore either contained in $u^{-1}(\alpha)$ or disjoint from it.

    Suppose that there are $k$ connected components of the critical locus that are contained in $u^{-1}(\alpha)$.
    If $k = 0$, then $u^{-1}(\alpha)$ contains no critical points of $f$, so we can choose $C > C_1 > 0$ such that $\|\nabla v\| > C_1$ on $M\setminus K$ and on $u^{-1}(\alpha)\cap K$, hence on $u^{-1}(\alpha)$.
    Then by Lemma~\ref{lem:integral-curve}, $v\big(u^{-1}(\alpha)\big) = \mathbb{R}$.
    Otherwise, let $k>0$ and let $S_1, \ldots, S_k$ be the connected components of the critical locus that are contained in $u^{-1}(\alpha)$, arranged so that
    \begin{equation}
        v(S_1) \leq \cdots \leq v(S_k).
    \end{equation}

    Since $f$ is a nonconstant holomorphic function, it is an open map.
    Thus there exists $\epsilon > 0$ such that for each $i$,
    \begin{equation}
        D_\epsilon\big(f(S_i)\big) \subset f(M),
    \end{equation}
    and for $i \neq j$ and $v(S_i) \neq v(S_j)$,
    \begin{equation}
        D_\epsilon\big(f(S_i)\big) \cap D_\epsilon\big(f(S_j)\big) = \emptyset.
    \end{equation}
    Hence
    \begin{equation}
        \big(v(S_i) - \epsilon, \, v(S_i) + \epsilon\big) \subset v\big(u^{-1}(\alpha)\big).
    \end{equation}

    Choose an open neighborhood $U_\alpha \subset M$ containing $S_1, \ldots, S_k$ such that
    \begin{equation}
        f(U_\alpha) \subset \bigcup_{i=1}^{k} D_{\epsilon'}\big(f(S_i)\big)
    \end{equation}
    for some $0 < \epsilon' < \epsilon$.

    Then
    \begin{equation}
        (K \setminus U_\alpha) \cap u^{-1}(\alpha)
    \end{equation}
    is compact, and there exists $\widetilde{C}_\alpha > 0$ such that
    \begin{equation}
        \|\nabla v\| > \widetilde{C}_\alpha \quad \text{on } u^{-1}(\alpha) \setminus U_\alpha.
    \end{equation}

    Choose $p_1 \in u^{-1}(\alpha) \setminus U_\alpha$ such that
    \begin{equation}
        v(p_1) \in \big(v(S_1) - \epsilon, \, v(S_1) - \epsilon'\big).
    \end{equation}

    Consider the maximal forward integral curve $\gamma_0$ of $-\nabla v$ starting at $p_1$.
    Since $v$ strictly decreases along this curve, the curve cannot intersect $U_{\alpha}$.
    Let $\Gamma_0$ be the image of $\gamma_0$.
    By Lemma~\ref{lem:integral-curve},
    \begin{equation}
        v(\Gamma_0) = (-\infty, v(p_1)].
    \end{equation}

    Set $j_1 = 1$.
    Choose $q_1 \in u^{-1}(\alpha) \setminus U_{\alpha}$ such that
    \begin{equation}
        v(q_1) \in \bigl(v(S_{j_1}) + \epsilon', \ v(S_{j_1}) + \epsilon\bigr).
    \end{equation}

    Consider the maximal forward integral curve $\gamma_1$ of $\nabla v$ starting at $q_{j_1}$.
    If $\gamma_1$ does not intersect $U_{\alpha}$, then by Lemma~\ref{lem:integral-curve}
    \begin{equation}
        v(\gamma_1) = [v(q_{j_1}), +\infty),
    \end{equation}
    and the argument terminates.

    Otherwise, there exist $t_1$ and $j_2 > j_1$ such that
    \begin{equation}
        p_{j_2} := \gamma_1(t_1) \in U_{\alpha}, \quad v(p_{j_2}) \in \bigl(v(S_{j_2}) - \epsilon', \ v(S_{j_2}) + \epsilon'\bigr).
    \end{equation}
    Let
    \begin{equation}
        \Gamma_1 = \gamma_1([0, t_1]).
    \end{equation}

    Repeating the same procedure, we obtain indices
    \begin{equation}
        1 = j_1 < j_2 < \cdots < j_m \le k
    \end{equation}
    for some $1 \le m \le k$, points $p_{j_2}, \ldots, p_{j_m} \in U_{\alpha}$ such that
    \begin{equation}
        v(p_{j_i}) \in \bigl(v(S_{j_i}) - \epsilon', \ v(S_{j_i}) + \epsilon'\bigr),
    \end{equation}
    points $q_{j_1}, \ldots, q_{j_m} \in u^{-1}(\alpha) \setminus U_{\alpha}$ such that
    \begin{equation}
        v(q_{j_i}) \in \bigl(v(S_{j_i}) + \epsilon', \ v(S_{j_i}) + \epsilon\bigr),
    \end{equation}
    and curve segments $\Gamma_1, \ldots, \Gamma_m$, where $\Gamma_k$ connects $q_{j_k}$ to $p_{j_{k+1}}$ for $k = 1, \ldots, m-1$ and $\Gamma_m$ is the image of a maximal forward integral curve of $\nabla v$.
    By Lemma~\ref{lem:integral-curve},
    \begin{equation}
        v(\Gamma_m) = [v(q_{j_m}), +\infty).
    \end{equation}

    The sets
    \begin{equation}
        v(\Gamma_0), \ v(\Gamma_1), \ldots, \ v(\Gamma_m)
    \end{equation}
    are pairwise disjoint closed intervals whose complement in $\mathbb{R}$ is
    \begin{equation}
        \bigl(v(p_{j_1}), v(q_{j_1})\bigr) \cup \cdots \cup \bigl(v(p_{j_m}), v(q_{j_m})\bigr).
    \end{equation}

    Each of these complementary intervals is contained in one of the open intervals
    \begin{equation}
        \bigl(v(S_i) - \epsilon, \ v(S_i) + \epsilon\bigr),
    \end{equation}
    which in turn lie in $v\bigl(u^{-1}(\alpha)\bigr)$.
    Therefore
    \begin{equation}
        v\bigl(u^{-1}(\alpha)\bigr) = \mathbb{R}.
    \end{equation}

    Hence $f(M) = \mathbb{C}$.
\end{proof}

As an application of the range theorem, we record the following corollary concerning holomorphic embeddings and polynomial superpotentials.

\begin{corollary}\label{cor:embedding}
    Let $(M,f)$ be a K\"ahlerian LG model, and let $\psi: N\hookrightarrow M$ be a holomorphic embedding.
    If the induced K\"ahlerian LG model $(N, \psi^* f)$ is fundamentally tame, then $f(M)=\mathbb{C}$.
    In particular, the following two cases hold:
    \begin{enumerate}
        \item If $f$ is a nonconstant polynomial defined on the Euclidean space $\mathbb{C}^N$, then $f(\mathbb{C}^N)=\mathbb{C}$.
        \item If $f$ is a Laurent polynomial defined on $\bigl((\mathbb{C}^*)^N, \sum_{i=1}^N |dz_i|^2/|z_i|^2\bigr)$ having one monomial with positive degree and one monomial with negative degree, then $f\bigl((\mathbb{C}^*)^N\bigr)=\mathbb{C}$.
    \end{enumerate}
\end{corollary}

\begin{proof}
    Since $(N, \psi^* f)$ is fundamentally tame, the range theorem (Theorem~\ref{thm:range}) applied to this LG model gives $\psi^*f(N)=\mathbb{C}$, and hence $f(M)=\mathbb{C}$.
    For the polynomial case, choose $v\in\mathbb{C}^N$ such that $z\mapsto f(zv)$ is nonconstant.
    This one-variable polynomial is surjective by the fundamental theorem of algebra, so $f(\mathbb{C}^N)=\mathbb{C}$.
    For the Laurent case, choose $\alpha\in(\mathbb{C}^*)^N$ generically so that $h(z)=f(\alpha_1z,\ldots,\alpha_Nz)$ retains a positive and a negative power of $z$.
    Such a choice is possible because the coefficient of each total degree is a nonzero Laurent polynomial in $\alpha$.
    If $-r<0$ is the lowest exponent of $h$, then for every $c\in\mathbb{C}$, $z^r(h(z)-c)$ is a nonconstant polynomial with nonzero constant term.
    It therefore has a nonzero root, which gives $h(z)=c$ and proves the second case.
\end{proof}

The following class of quiver superpotentials gives examples satisfying the hypothesis of Corollary~\ref{cor:embedding}~(2).

\begin{example}\label{exa:quiver}
    Let $Q=(Q_0,Q_1)$ be a finite directed quiver with $|Q_0|\ge2$, vertex set $Q_0$, and arrow set $Q_1$.
    Attach a nonzero variable $y_v\in\mathbb{C}^*$ to every vertex $v\in Q_0$, and to an arrow $a:u\to v$ associate the character
    \begin{equation*}
        \chi_a=\frac{y_v}{y_u}.
    \end{equation*}
    The associated quiver superpotential is the Laurent polynomial
    \begin{equation*}
        W_Q=\sum_{a\in Q_1}c_a\,\chi_a
        =\sum_{a\in Q_1}c_a\,\frac{y_{h(a)}}{y_{t(a)}},
        \qquad c_a\in\mathbb{C}^*,
    \end{equation*}
    where $t(a)$ and $h(a)$ denote the tail and the head of the arrow $a$.
    Since the simultaneous rescaling $y_v\mapsto\lambda y_v$ leaves every ratio $y_{h(a)}/y_{t(a)}$ invariant, the natural habitat of $W_Q$ is the algebraic torus
    \begin{equation*}
        T_Q=(\mathbb{C}^*)^{Q_0}/\mathbb{C}^*\simeq(\mathbb{C}^*)^{|Q_0|-1},
    \end{equation*}
    where one fixes a base vertex $v_0\in Q_0$ with $y_{v_0}=1$, endowed with the standard logarithmic metric $\sum_i |dy_i|^2/|y_i|^2$.

    The exponent of the arrow $a:u\to v$ is the vector $m_a=e_v-e_u$ in the character lattice $\{m\in\mathbb{Z}^{Q_0}:\sum_{v\in Q_0}m_v=0\}$.
    After combining equal characters, the Newton polytope of $W_Q$ is
    \begin{equation*}
        \Delta_Q=\operatorname{Conv}\{m:c_m\ne0\},\qquad c_m=\sum_{a\in Q_1:\,m_a=m}c_a.
    \end{equation*}
    Assume that $\Delta_Q$ is full-dimensional and that $0\in\operatorname{Int}(\Delta_Q)$.
    Then the actual support of $W_Q$ contains exponents of both positive and negative total degree in the chosen torus coordinates, so the hypothesis of Corollary~\ref{cor:embedding}~(2) is satisfied, and we conclude that
    \begin{equation*}
        W_Q(T_Q)=\mathbb{C}.
    \end{equation*}

    For $n\ge2$ and $q\in\mathbb{C}^*$, a concrete instance is the Hori--Vafa mirror of $X=\mathbb{P}^{n-1}$, namely
    \begin{equation*}
        W=x_1+\cdots+x_{n-1}+\frac{q}{x_1\cdots x_{n-1}}
        \qquad\text{on }(\mathbb{C}^*)^{n-1},
    \end{equation*}
    whose monomials $x_i$ have positive degree and whose monomial $\frac{q}{x_1\cdots x_{n-1}}$ has negative degree; hence $W\bigl((\mathbb{C}^*)^{n-1}\bigr)=\mathbb{C}$ by Corollary~\ref{cor:embedding}~(2).
\end{example}

\section{Fibration Theorem}\label{sec:fibration}

In this section we establish a fibration theorem for the superpotential $f$.
It will be used both in the computation of relative homology and in the comparison between fibers and sublevel sets.

\begin{proposition}\label{prop:fibration}
    Let $(M,f)$ be a fundamentally tame LG model.
    Denote by $D$ the set of critical values of $f$.
    Then
    \begin{equation}
        f|_{f^{-1}(f(M)\setminus D)} : f^{-1}(f(M) \setminus D) \to f(M) \setminus D
    \end{equation}
    is a locally trivial fibration.
\end{proposition}

The argument follows the strategy of the Ehresmann fibration theorem and is inspired by \cite[Lemma 4.3]{Broughton1988}.
Since $f$ is a submersion over $f(M) \setminus D$, a trivialization is obtained by lifting a vector field from the base and constructing an isotopy between fibers along its flow.
As $f$ is not assumed to be proper, completeness of the lifted flow is not automatic; we therefore construct the lifted vector field explicitly, and use the fundamentally tame condition together with the completeness of $M$ to ensure that its flow exists globally.

\begin{proof}
    Let $U \subset f(M)$ be a convex subset with $\overline{U} \subset f(M) \setminus D$.
    Since $f(M) \setminus D$ is open in $\mathbb{C}$, we may choose an open neighborhood $V$ of $\overline{U}$ with $\overline{V} \subset f(M) \setminus D$.
    Then $f$ is a submersion on $f^{-1}(\overline{V})$.

    We look for a vector field $\tilde{X}$ on $f^{-1}(V)$ lifting the constant vector field on the base, that is, satisfying
    \begin{equation}
        \mathrm{d}f(\tilde{X}) = 1.
    \end{equation}
    Since $f$ is holomorphic, $\partial f$ is a $(1,0)$-form, and its gradient
    \begin{equation}
        \nabla f = g^{i\bar{j}} \frac{\partial f}{\partial z^i} \frac{\partial}{\partial \bar{z}^j}
    \end{equation}
    is a $(0,1)$-vector field.
    Define
    \begin{equation}
        \tilde{X} = \frac{1}{\|\nabla f\|^2} \overline{\nabla f}.
    \end{equation}
    Then $\tilde{X}$ is a smooth $(1,0)$-vector field, and a direct computation gives
    \begin{equation}
        \tilde{X} f = \frac{1}{\|\nabla f\|^2} \overline{\nabla f}(f) = \frac{1}{\|\nabla f\|^2} \|\nabla f\|^2 = 1.
    \end{equation}
    Hence $\partial f(\tilde{X}) = 1$, so that $\tilde{X}$ lifts the constant unit vector field on the base.

    Because $f$ is fundamentally tame and has no critical point on $f^{-1}(\overline{V})$, the function $\|\nabla f\|$ admits a positive lower bound on $f^{-1}(\overline{V})$.
    In particular $\tilde{X}$ is bounded.

    Choose a smooth function $\mu$ on $M$ such that
    \begin{equation}
        \mu \equiv 1 \ \text{ on } \ f^{-1}(U), \qquad \mu \equiv 0 \ \text{ on } \ M \setminus f^{-1}(V), \qquad 0 \le \mu \le 1,
    \end{equation}
    and define a global smooth vector field on $M$ by
    \begin{equation}
        X(p) =
        \begin{cases}
            \mu(p)\tilde{X}(p), & p \in f^{-1}(V), \\
            0, & p \notin f^{-1}(V).
        \end{cases}
    \end{equation}
    Then $X$ is bounded.

    Fix $c_0 \in U$ and, for each $c \in U$, set
    \begin{equation}
        X_c = (c_0 - c)X,
    \end{equation}
    denoting the global flow of $2 \operatorname{Re} X_c$ by $\phi_c(p,t)$.
    Since $X_c$ is bounded and $M$ is complete, this flow is complete.

    Then for $p \in f^{-1}(c)$ we have
    \begin{equation}
        \frac{\mathrm{d}}{\mathrm{d}t} f(\phi_c(p,t)) = (c_0 - c)\mu(\phi_c(p,t)).
    \end{equation}
    Integrating yields
    \begin{equation}
        f(\phi_c(p,t)) = c + \int_0^t (c_0 - c)\mu(\phi_c(p,s)) \, \mathrm{d}s.
    \end{equation}

    Since $0 \le \mu \le 1$, for $t \in [0,1]$ we may write
    \begin{equation}
        f(\phi_c(p,t)) = (1 - \theta(t))c + \theta(t)c_0
    \end{equation}
    for some $\theta(t) \in [0,1]$; hence $f(\phi_c(p,t))$ lies on the segment joining $c$ and $c_0$.
    Because $U$ is convex, this segment is contained in $U$, so $\mu = 1$ along the trajectory for $t \in [0,1]$.
    Therefore
    \begin{equation}
        f(\phi_c(p,t)) = (1 - t)c + tc_0.
    \end{equation}

    It follows that
    \begin{equation}
        \phi_c(f^{-1}(c), t) \subset f^{-1}\bigl((1 - t)c + tc_0\bigr).
    \end{equation}
    Since the flow is invertible, with inverse given by $t \mapsto -t$, the reverse inclusion also holds.
    Hence
    \begin{equation}
        \phi_c(f^{-1}(c), t) = f^{-1}\bigl((1 - t)c + tc_0\bigr),
    \end{equation}
    and the induced map
    \begin{equation}
        \tilde{\phi}_{c,t} : f^{-1}(c) \to f^{-1}\bigl((1 - t)c + tc_0\bigr)
    \end{equation}
    is a diffeomorphism.

    Define
    \begin{equation}
        H(p,t) = \phi_{f(p)}(p,t), \qquad p \in f^{-1}(U), \ t \in [0,1].
    \end{equation}
    Then
    \begin{equation}
        \frac{\partial H}{\partial t} = 2\operatorname{Re}((c_0 - f(p))X)(H(p,t)), \qquad H(p,0) = p.
    \end{equation}
    Thus $H$ is smooth and
    \begin{equation}
        H(f^{-1}(U), 1) = f^{-1}(c_0).
    \end{equation}
    Hence $H$ defines a deformation retraction of $f^{-1}(U)$ onto $f^{-1}(c_0)$.

    Finally define
    \begin{equation}
        \tilde{H} : f^{-1}(U) \to f^{-1}(c_0) \times U, \qquad \tilde{H}(p) = (H(p,1), f(p)).
    \end{equation}
    Its inverse is obtained by flowing backwards along the corresponding segment in $U$.
    Therefore $\tilde{H}$ is a diffeomorphism, and $f|_{f^{-1}(U)}$ is a trivial fibration over $U$.
\end{proof}

\section{Relative homology of smooth fibers}\label{sec:relative}

In this section we compute the relative homology of the regular fibers of $f$.
The resulting description of $H_*(M, f^{-1}(c))$ will be used in the next section to construct the relative cohomology bundle over the deformation parameter space.

\begin{theorem}\label{thm:main}
    Let $(M,f)$ be a fundamentally tame LG model.
    Assume that $f$ has finitely many critical points and set $\mu = \mu_f$.
    Then for any regular value $c$ of $f$ the relative homology groups $H_*(M, f^{-1}(c))$ are concentrated in degree $n$, and
    \begin{equation}
        H_n(M, f^{-1}(c)) \cong \mathbb{Z}^{\mu}.
    \end{equation}
\end{theorem}

We shall use the following lemma, which is a direct consequence of Lemma B.3 in \cite{milnor1968singular} applied to $\partial f$.

\begin{lemma}\label{lem:perturbation}
    Let $V \subset \mathbb{C}^n$ be an open neighborhood of the origin.
    Let $f : V \to \mathbb{C}$ be a holomorphic function such that the origin is the only critical point of $f$ in $V$.
    Let $U \Subset V$ be an open neighborhood of the origin.
    Then there exists $\epsilon > 0$ such that for almost all tuples $a = (a_1, \ldots, a_n) \in \mathbb{C}^n$ satisfying $\|a\| < \epsilon$, the perturbed function
    \begin{equation}
        \tilde{f}(z) = f(z) + a_1 z_1 + \cdots + a_n z_n
    \end{equation}
    is a Morse function in $U$.
    Moreover, the number of critical points of $\tilde{f}$ equals the Milnor number $\mu_f$.
\end{lemma}

\begin{proof}[Proof of Theorem~\ref{thm:main}]
    We follow the procedure in Section 4 of \cite{Broughton1988} to construct a CW approximation for $(M, f^{-1}(c))$.

    \medskip
    \noindent\emph{Step 1.}
    Let $p_1, \ldots, p_k$ be all critical points of $f$.
    Choose pairwise disjoint coordinate charts $(V_i, (z^i_1, \ldots, z^i_n))$ centered at $p_i$ such that $\overline{V_i}$ is compact and $\overline{V_i} \cap f^{-1}(c) = \varnothing$, and open neighborhoods $U_i$ of $p_i$ such that $\overline{U_i} \subset V_i$.
    Choose smooth functions $\phi^i$ supported in $V_i$ such that $\phi^i = 1$ on $U_i$.
    For each $i$, let $a^i = (a^i_1, \ldots, a^i_n) \in \mathbb{C}^n$ be a tuple of complex numbers.
    On $V_i$ define
    \begin{equation}
        Z_i = (z^i_1 \phi^i, \ldots, z^i_n \phi^i),
    \end{equation}
    and
    \begin{equation}
        A^i_{a^i} = \biggl( \sum_{j=1}^{n} a^i_j z^i_j \biggr) \phi^i = \langle a^i, Z_i \rangle, \qquad A_a = \sum_{i=1}^{k} A^i_{a^i},
    \end{equation}
    where each $A^i_{a^i}$ is understood to be zero outside $V_i$.
    We then define the perturbed function $f_a$ by
    \begin{equation}
        f_a = f + A_a.
    \end{equation}

    Define
    \begin{equation}
        F_a = |f_a - c|^2, \qquad F = |f - c|^2.
    \end{equation}
    Clearly $F_a$ coincides with $F$ outside $\bigcup_{i=1}^{k} V_i$.
    Let
    \begin{equation*}
        \eta = \min_{\bigcup_i \overline{V_i}} |f - c| > 0.
    \end{equation*}
    We can restrict the norm of $a$ such that $\|A_a\|_\infty < \eta/2$, thus $F_a \neq 0$ on $\bigcup_{i=1}^{k} \overline{V_i}$.
    Then $F_a^{-1}(0) = F^{-1}(0) = f^{-1}(c)$.
    Let $\Delta$ be the open neighborhood of $f^{-1}(c)$ given by the inequality $F_a < \delta^2$ where $0 < \delta < \eta/2$, thus $\overline{\Delta} \subset M \setminus \bigcup_{i=1}^{k} \overline{V_i}$.

    We first prove that $F_a$ is fundamentally tame on $M \setminus \Delta$ and that, if each $\|a^i\|$ is sufficiently small, $F_a$ has no critical point on $M \setminus (\Delta \cup \bigcup_{i=1}^{k} U_i)$.
    Since $F$ is real-valued, $\mathrm{d}F = 0$ is equivalent to $\partial F = 0$, and we compute
    \begin{equation}
        \partial F = \overline{(f-c)}\,\partial f.
    \end{equation}
    Since $f$ is fundamentally tame, there exists a compact subset $K \subset M$ and a positive constant $C$ such that $\|\partial f\| = \|\nabla f\| > C$ on $M \setminus K$; we may choose $K$ so that $\bigcup_{i=1}^{k} \overline{V_i} \subset K$.
    Then on the set $M \setminus (K \cup \Delta)$ we have
    \begin{equation}
        \|\partial F\| = |f-c|\,\|\partial f\| > \delta C.
    \end{equation}
    Since $F_a$ coincides with $F$ outside $K$, the inequality above holds for $\|\partial F_a\|$ as well; hence $F_a$ is fundamentally tame on $M \setminus \Delta$.

    We now consider the behavior of $F_a$ on $M \setminus \bigcup_{i=1}^{k} U_i$.
    Since $\partial f \neq 0$ and $f \neq c$ on this set, $F$ has no critical point there.
    Since $F_a$ coincides with $F$ outside $\bigcup_{i=1}^{k} \overline{V_i}$, it suffices to show that $F_a$ has no critical point on $\overline{V_i \setminus U_i}$ for each $i$.
    We argue by estimating $\|\mathrm{d}F_a - \mathrm{d}F\|$ on these sets.
    Observe
    \begin{equation}
        F_a = F + 2\operatorname{Re}\bigl( (f-c)\overline{A_a} \bigr) + |A_a|^2.
    \end{equation}
    Hence
    \begin{equation}
        \mathrm{d}F_a - \mathrm{d}F = \mathrm{d}\bigl( 2\operatorname{Re}((f-c)\overline{A_a}) \bigr) + \mathrm{d}(|A_a|^2).
    \end{equation}
    We obtain
    \begin{equation}
        \begin{aligned}
            \|\mathrm{d}F_a - \mathrm{d}F\| &= \bigl\| \mathrm{d}\bigl( 2\operatorname{Re}((f-c)\overline{A_a}) \bigr) + \mathrm{d}(|A_a|^2) \bigr\| \\
            &= \bigl\| 2\operatorname{Re}(\overline{A_a}\,\mathrm{d}f + (f-c)\,\mathrm{d}\overline{A_a}) + 2\operatorname{Re}(\overline{A_a}\,\mathrm{d}A_a) \bigr\| \\
            &\le 2 \bigl( |A_a|\,\|\mathrm{d}f\| + |f-c|\,\|\mathrm{d}A_a\| + |A_a|\,\|\mathrm{d}A_a\| \bigr) \\
            &= 2 \bigl( |\langle a^i, Z_i \rangle|\,\|\mathrm{d}f\| + |f-c|\,\|\langle a^i, \mathrm{d}Z_i \rangle\| \\
            &\qquad {}+ |\langle a^i, Z_i \rangle|\,\|\langle a^i, \mathrm{d}Z_i \rangle\| \bigr) \\
            &\le 2 \bigl( \|a^i\|\,\|Z_i\|\,\|\mathrm{d}f\| + |f-c|\,\|a^i\|\,\|\mathrm{d}Z_i\| \\
            &\qquad {}+ \|a^i\|\,\|Z_i\|\,\|a^i\|\,\|\mathrm{d}Z_i\| \bigr).
        \end{aligned}
    \end{equation}

    On the compact set $\overline{V_i \setminus U_i}$ the following functions and forms are bounded:
    \begin{equation}
        Z_i, \quad \mathrm{d}Z_i, \quad f-c, \quad \partial f.
    \end{equation}
    Hence there exists a constant $K_i > 0$ such that
    \begin{equation}
        \|\mathrm{d}F_a - \mathrm{d}F\| \le K_i \|a^i\|
    \end{equation}
    for $\|a^i\|$ sufficiently small.

    Moreover, since $\partial f$ does not vanish on $\overline{V_i \setminus U_i}$, there exists $m_i > 0$ such that
    \begin{equation}
        \|\mathrm{d}F\| = 2|f-c|\,\|\partial f\| \ge m_i
    \end{equation}
    on $\overline{V_i \setminus U_i}$.

    Choose $\epsilon > 0$ such that
    \begin{equation}
        K_i \epsilon \le \frac{m_i}{2}
    \end{equation}
    for every $i$.
    Then whenever $\|a^i\| < \epsilon$,
    \begin{equation}
        \|\mathrm{d}F_a - \mathrm{d}F\| \le \frac{1}{2}\|\mathrm{d}F\|
    \end{equation}
    on $\overline{V_i \setminus U_i}$.

    Second, by Lemma~\ref{lem:perturbation}, we may shrink $\epsilon$ so that for almost all $a^i$ with $\|a^i\| < \epsilon$, the function $f_a$ is Morse on $U_i$ and has exactly $\mu_f(p_i)$ critical points in $U_i$.
    We fix such an $a^i$ for each $i$, chosen so that $F_a$ takes distinct values at distinct critical points; this will be used in Step 2.
    Since $f_a$ is holomorphic on $U_i$ while $F_a$ is real-valued, the critical points of $F_a$ in $U_i$ are determined by
    \begin{equation}
        \partial F_a = \overline{(f_a - c)}\,\partial f_a = 0.
    \end{equation}
    Thus they coincide with the critical points of $f_a$ in $U_i$, and their number equals $\mu_f(p_i)$.

    Finally, we compute the Morse index of each critical point of $F_a$.
    Since $f_a$ is holomorphic on $U_i$, the Hessian of $F_a$ in $(z^i, \bar{z}^{\,i})$ coordinates has the form
    \begin{equation}
        \operatorname{Hess} F_a =
        \begin{pmatrix}
            \overline{(f_a - c)}\,\partial^2 f_a & (\partial f_a)(\overline{\partial f_a})^T \\
            (\overline{\partial f_a})(\partial f_a)^T & (f_a - c)\,\overline{\partial^2 f_a}
        \end{pmatrix}.
    \end{equation}
    Since $f_a$ differs from $f$ by linear terms on $U_i$, their second derivatives coincide:
    \begin{equation}
        \partial^2 f_a = \partial^2 f.
    \end{equation}
    At a critical point $p \in U_i$ we have $\partial f_a(p) = 0$, so the off-diagonal blocks of the Hessian vanish.
    Denote the nonzero constant $\overline{(f_a(p)-c)}$ by $\lambda_p$.
    Then the Hessian of $F_a$ at $p$ has the form
    \begin{equation}
        \operatorname{Hess} F_a(p) =
        \begin{pmatrix}
            \lambda_p\,\partial^2 f(p) & 0 \\
            0 & \overline{\lambda_p\,\partial^2 f(p)}
        \end{pmatrix}.
    \end{equation}
    It equals the Hessian at $p$ of the function $\lambda_p f_a + \overline{\lambda_p f_a} = 2\operatorname{Re}(\lambda_p f_a)$.
    Since $\lambda_p \neq 0$, the functions $F_a$ and $2\operatorname{Re}(f_a)$ have the same Morse index at $p$.
    By standard complex Morse theory, the real part of a holomorphic function on a complex $n$-manifold has Morse index $n$ at a nondegenerate critical point.
    Hence the Morse index of $F_a$ at $p$ equals $n$.

    \medskip
    \noindent\emph{Step 2.} We construct a CW approximation for $(M, f^{-1}(c))$ by a Morse-theoretic argument in the spirit of Milnor \cite{milnor1963morse} and Broughton \cite{Broughton1988}.

    Recall that, by the definition of $\Delta$, $\overline{\Delta} = f_a^{-1}(\overline{D_\delta(c)}) = f^{-1}(\overline{D_\delta(c)})$, where $D_\delta(c)$ denotes the open disk of radius $\delta$ centered at $c$ in $\mathbb{C}$.
    By Proposition~\ref{prop:fibration}, $f^{-1}(c)$ is a deformation retract of $\overline{\Delta}$; it therefore suffices to understand the homotopy type of $(M, \overline{\Delta})$.

    Let
    \begin{equation}
        M^r = F_a^{-1}[0, r], \qquad M^{s,r} = F_a^{-1}[s, r].
    \end{equation}
    Clearly $M^{\delta^2} = \overline{\Delta}$.

    As in Theorems 3.1 and 3.2 of \cite{milnor1963morse}, we claim:
    \begin{enumerate}
        \item If $r \ge s \ge \delta^2$ and $F_a$ has no critical point in $M^{s,r}$, then $M^s$ is a deformation retract of $M^r$.
        \item If $p$ is a critical point of index $k$ with $F_a(p) = s$, then for sufficiently small $\epsilon > 0$, $M^{s+\epsilon}$ has the homotopy type of $M^{s-\epsilon}$ with a $k$-cell attached.
    \end{enumerate}
    Then induction on the critical values of $F_a$ shows that $M$ has the homotopy type of $\overline{\Delta}$ with $\mu$ $n$-cells attached.

    In Milnor's argument the sets $M^{s,r}$ were assumed compact.
    As observed by Broughton, the argument remains valid under the weaker assumptions that
    \begin{enumerate}
        \item each $M^{s,r}$ is complete with respect to the given Riemannian metric;
        \item $\|\nabla F_a\|$ is bounded away from $0$ outside a compact neighborhood of the critical points in $M^{s,r}$.
    \end{enumerate}
    In the present situation $M$ is complete and $F_a$ is fundamentally tame on $M \setminus \Delta$, so the two conditions above are satisfied for every $M^{s,r}$ with $s \ge \delta^2$.
    Therefore the proofs of Theorems 3.1 and 3.2 in \cite{milnor1963morse} apply without modification, and our claims hold.

    We conclude that $M$ has the homotopy type of $f^{-1}(c)$ with $\mu$ $n$-cells attached.
    Consequently,
    \begin{equation}
        H_k(M, f^{-1}(c)) = 0 \quad \text{for } k \neq n,
    \end{equation}
    and
    \begin{equation}
        H_n(M, f^{-1}(c)) \cong \mathbb{Z}^{\mu}.
    \end{equation}
\end{proof}

\begin{remark}\label{rem:broughton}
    When $M = \mathbb{C}^{n+1}$ and $f$ is a tame polynomial in the sense of Broughton, the statement above reduces to a special case of \cite[Theorem 1.2]{Broughton1988}, where only regular fibers are considered.
    Since $\mathbb{C}^{n+1}$ is contractible, the homology of a regular fiber can be recovered from the relative CW approximation of the pair $(M, f^{-1}(c))$ obtained from the Morse-theoretic construction.

    In the present setting, however, $M$ is a complete K\"ahler manifold and need not be contractible.
    What one obtains is therefore the relative homology of the pair $(M, f^{-1}(c))$ rather than the homology of the fiber itself.
\end{remark}

\section{Construction of Cohomological Milnor Fibrations}\label{sec:construction}

In this section we study deformations of the superpotential $f$ and the fibration structure they induce.
Using a fibration theorem analogous to Proposition~\ref{prop:fibration}, we show that nearby fibers are related by isotopies in $M$ depending smoothly on the deformation parameter.
It follows that the relative cohomology groups $H^n(M, f_t^{-1}(c))$ form a smooth vector bundle over the parameter space, equipped with a natural flat connection.
To relate this construction to the Lefschetz thimbles of Picard--Lefschetz theory, we also identify the relative homotopy type of the fibers with that of the corresponding sublevel sets at infinity.

Let $S \subset \mathbb{C}^m$ be a bounded open neighborhood of $0$.
We consider deformations of $f$ of the form
\begin{equation}
    F(z,t) = f(z) + \sum_{i=1}^{m} t^i g_i(z), \quad (z,t) \in M \times S,
\end{equation}
where the $g_i$ are holomorphic functions on $M$ and $g_1 = 1$.

For each fixed $t \in S$, denote
\begin{equation}
    f_t(z) = F(z,t).
\end{equation}

\begin{theorem}\label{thm:deformation}
    Suppose that there exist a compact subset $K \subset M$ and constants $C_i > 0$ for $i = 1, \ldots, m$ such that for every $t \in S$,
    \begin{equation}
        \|\nabla f_t(z)\| > C_i |g_i(z)| \quad \text{for all } z \in M \setminus K.
    \end{equation}
    Then for $c \in \mathbb{C}$ with $|\operatorname{Re} c|$ sufficiently large, the level set
    \begin{equation}
        E_c = \{(z,t) \in M \times S \mid F(z,t) = c\}
    \end{equation}
    is a smooth hypersurface of $M \times S$, and the projection
    \begin{equation}
        \pi : E_c \to S
    \end{equation}
    is a locally trivial smooth fibration.
\end{theorem}

\begin{proof}
    First, we prove that $E_c$ is a smooth hypersurface of $M \times S$.

    Taking $i = 1$ in the assumption, we obtain $\|\nabla f_t(z)\| > C_1$ for all $z \in M \setminus K$ and all $t \in S$.
    Hence all critical points of $f_t$ are contained in $K$, uniformly for $t \in S$.
    Moreover, since $S$ is bounded, the values of $f_t$ on $K$ are uniformly bounded for $t \in S$; choose $R > 0$ such that
    \begin{equation}
        |\operatorname{Re} f_t(z)| \le R \quad \text{for all } z \in K \text{ and all } t \in S.
    \end{equation}
    Then for any $c \in \mathbb{C}$ with $|\operatorname{Re} c| > R$, $E_c\cap(K\times S)=\varnothing$.
    Thus for every $(z,t) \in E_c$ we have
    \begin{equation}
        \|\nabla_z F(z,t)\| = \|\nabla f_t(z)\| \neq 0.
    \end{equation}
    Here $\nabla_z F$ denotes the gradient of $F$ with respect to the $z$-variables only.
    Hence $\mathrm{d}F(z,t) \neq 0$ on $E_c$, and $E_c$ is a smooth hypersurface of $M \times S$.

    Next, we prove that $\pi : E_c \to S$ is a locally trivial smooth fibration.

    Fix a point
    \begin{equation}
        t_0 = (t_0^1, \ldots, t_0^m) \in S.
    \end{equation}
    Since $S \subset \mathbb{C}^m$ is open, there exist discs
    \begin{equation}
        D_i \subset \mathbb{C}, \quad 1 \le i \le m,
    \end{equation}
    each centered at $t_0^i$, such that
    \begin{equation}
        D_1 \times \cdots \times D_m \subset S.
    \end{equation}
    For $1 \le k \le m$, denote
    \begin{equation}
        U_k = D_1 \times \cdots \times D_k \times \{t_0^{k+1}\} \times \cdots \times \{t_0^m\}.
    \end{equation}
    It suffices to prove that
    \begin{equation}
        \pi^{-1}(U_m) \to U_m
    \end{equation}
    is trivial over $U_m$.
    We argue by induction on $m$; the case $m = 0$ is trivial.
    Assume that for $m - 1$ the projection is trivial over $U_{m-1}$ and that there exists a smooth deformation retraction of $\pi^{-1}(U_{m-1})$ onto $\pi^{-1}(U_0) = \pi^{-1}(t_0)$ whose restriction to $\pi^{-1}(t^1, \ldots, t^{m-1}, t_0^m)$ is an isotopy for every $(t^1, \ldots, t^{m-1}, t_0^m) \in U_{m-1}$.

    On $E_c$, define the vector field
    \begin{equation}
        X_m = \frac{\partial}{\partial t^m} - \frac{g_m(z)}{\|\nabla_z F(z,t)\|^2}\, \overline{\nabla_z F(z,t)}.
    \end{equation}
    We compute
    \begin{equation}
        \begin{aligned}
            X_m F(z,t)
            &= \frac{\partial F}{\partial t^m} - \frac{g_m(z)}{\|\nabla_z F(z,t)\|^2}\, \overline{\nabla_z F(z,t)}\, F \\
            &= g_m(z) - \frac{g_m(z)}{\|\nabla_z F(z,t)\|^2}\, \|\nabla_z F(z,t)\|^2 \\
            &= 0.
        \end{aligned}
    \end{equation}
    Hence $X_m$ is tangent to $E_c$.

    Moreover,
    \begin{equation}
        \pi_* X_m = \frac{\partial}{\partial t^m},
    \end{equation}
    so $X_m$ lifts the constant vector field $\partial/\partial t^m$.

    Take the standard Hermitian metric on $S$.
    By the assumption, for every $(z,t) \in E_c$ we have
    \begin{equation}
        \|X_m\| = \left\| \frac{\partial}{\partial t^m} - \frac{g_m(z)}{\|\nabla_z F(z,t)\|^2}\, \overline{\nabla_z F(z,t)} \right\| \le 1 + \frac{|g_m(z)|}{\|\nabla_z F(z,t)\|} < 1 + \frac{1}{C_m}.
    \end{equation}
    Hence $X_m$ is smooth and bounded on $E_c$.

    By an argument analogous to that in Proposition~\ref{prop:fibration}, there exists a smooth deformation retraction of $\pi^{-1}(U_m)$ onto $\pi^{-1}(U_{m-1})$ whose restriction to $\pi^{-1}(t^1, \ldots, t^{m-1}, t^m)$ is an isotopy for every $(t^1, \ldots, t^{m-1}, t^m) \in U_m$.
    By the induction hypothesis, there also exists a smooth deformation retraction of $\pi^{-1}(U_m)$ onto $\pi^{-1}(t_0)$ whose restriction to $\pi^{-1}(t^1, \ldots, t^m)$ is an isotopy for every $(t^1, \ldots, t^m) \in U_m$.
    This yields a trivialization of $\pi$ over $U_m$.
\end{proof}

We now clarify the relation between regular fibers and sublevel regions of $f$.

For $a > 0$, define
\begin{equation}
    M_{-}^{a} := \{p \in M : \operatorname{Re}(f(p)) \le -a\},
\end{equation}
\begin{equation}
    M_{+}^{a} := \{p \in M : \operatorname{Re}(f(p)) \ge a\}.
\end{equation}

\begin{proposition}\label{prop:homotopy}
    Let $(M, f)$ be a fundamentally tame LG model.
    There exist constants $R_{-} > 0$ and $R_{+} > 0$ such that:

    (1) For any $a, b \ge R_{-}$, the pairs $(M, M_{-}^{a})$ and $(M, M_{-}^{b})$ are homotopy equivalent.
    Moreover, for any $a \ge R_{-}$ and any $c \in \mathbb{C}$ with $\operatorname{Re}(c) \le -R_{-}$, the pair $(M, M_{-}^{a})$ is homotopy equivalent to $(M, f^{-1}(c))$; we denote their common homotopy type by $(M, f^{-\infty})$.

    (2) For any $a, b \ge R_{+}$, the pairs $(M, M_{+}^{a})$ and $(M, M_{+}^{b})$ are homotopy equivalent.
    Moreover, for any $a \ge R_{+}$ and any $c \in \mathbb{C}$ with $\operatorname{Re}(c) \ge R_{+}$, the pair $(M, M_{+}^{a})$ is homotopy equivalent to $(M, f^{-1}(c))$; we denote their common homotopy type by $(M, f^{+\infty})$.
\end{proposition}

\begin{proof}
    We prove the negative case; the positive case is analogous.

    For $a \ge 0$, set
    \begin{equation}
        \mathbb{C}_{-}^{a} := \{c \in \mathbb{C} : \operatorname{Re}(c) \le -a\}.
    \end{equation}
    Then $M_{-}^{a} = f^{-1}(\mathbb{C}_{-}^{a})$.

    Since $f$ is fundamentally tame, there exists $R_{-} > 0$ such that $\mathbb{C}_{-}^{R_{-}}$ contains no critical value of $f$.
    By Theorem~\ref{thm:range}, $\mathbb{C}_{-}^{R_{-}}$ is contained in the range of $f$; by Proposition~\ref{prop:fibration}, the restriction
    \begin{equation}
        f : M_{-}^{R_{-}} \to \mathbb{C}_{-}^{R_{-}}
    \end{equation}
    is a smooth locally trivial fibration.

    For $a \ge R_{-}$, the half-plane $\mathbb{C}_{-}^{a}$ is a deformation retract of $\mathbb{C}_{-}^{R_{-}}$.
    Since the fibration has no critical value over $\mathbb{C}_{-}^{R_{-}}$, this deformation retract lifts to a homotopy equivalence of pairs
    \begin{equation}
        (M, M_{-}^{a}) \simeq (M, M_{-}^{R_{-}}).
    \end{equation}
    Similarly, for any $c$ with $\operatorname{Re}(c) \le -R_{-}$, the singleton $\{c\}$ is a deformation retract of $\mathbb{C}_{-}^{R_{-}}$.
    Hence
    \begin{equation}
        (M, f^{-1}(c)) \simeq (M, M_{-}^{R_{-}}).
    \end{equation}
    This completes the proof.
\end{proof}

We now come to the main result of this section.
Combining the canonical homotopy type at infinity with the deformation fibration, we obtain the following global statement.

\begin{theorem}\label{thm:bundle}
    Assume the hypotheses of Theorem~\ref{thm:deformation} and suppose that $M$ is Stein.
    Let $(M, f_t^{\pm\infty})$ denote the canonical relative homotopy type defined in Proposition~\ref{prop:homotopy}.
    Then the relative cohomology groups
    \begin{equation}
        \mathcal{H}^{\pm} = \coprod_{t\in S} H^n(M, f_t^{\pm\infty}; \mathbb{C})
    \end{equation}
    form a smooth vector bundle over $S$ of rank $\mu$, endowed with a natural flat Gauss--Manin connection.

    Moreover, the total space
    \begin{equation}
        \mathcal{H}_{\mathrm{tot}} = \mathcal{H}^{+} \oplus \mathcal{H}^{-}
    \end{equation}
    is called the cohomological Milnor bundle associated with the deformation $F$.
\end{theorem}

\begin{proof}
    We prove the statement for $\mathcal{H}^{+}$; the case of $\mathcal{H}^{-}$ is analogous.
    By Theorem~\ref{thm:deformation}, for $c \in \mathbb{C}$ with $\operatorname{Re} c$ sufficiently large, the projection
    \begin{equation}
        \pi : E_c \to S
    \end{equation}
    is a locally trivial smooth fibration, where
    \begin{equation}
        E_c = \{(z,t) \in M \times S \mid F(z,t) = c\}.
    \end{equation}

    By the assumption, we may choose the constant $R_+$ of Proposition~\ref{prop:homotopy} uniformly in $t$, and then choose $c$ so that Proposition~\ref{prop:homotopy} applies to $f_t$ for all $t \in S$.
    Hence the pair $(M, f_t^{-1}(c))$ has the homotopy type $(M, f_t^{+\infty})$ for every $t \in S$.

    Fix $t_0 \in S$ and let $U \subset S$ be a contractible open neighborhood of $t_0$.
    By Theorem~\ref{thm:deformation}, there exists a smooth trivialization
    \begin{equation}
        \Phi : \pi^{-1}(U) \longrightarrow \pi^{-1}(t_0) \times U.
    \end{equation}
    For $t_1,t_2 \in U$, this induces a diffeomorphism
    \begin{equation}
        \Phi_{t_1,t_2} : \pi^{-1}(t_1) \longrightarrow \pi^{-1}(t_2),
    \end{equation}
    depending smoothly on $t_1, t_2$.

    Since the projection to $M$
    \begin{equation}
        p_t : \pi^{-1}(t) \longrightarrow f_t^{-1}(c)
    \end{equation}
    is a diffeomorphism for each $t$, we obtain a smooth family of diffeomorphisms
    \begin{equation}
        \Psi_{t_1,t_2} := p_{t_2} \circ \Phi_{t_1,t_2} \circ p_{t_1}^{-1} : f_{t_1}^{-1}(c) \longrightarrow f_{t_2}^{-1}(c).
    \end{equation}

    Because $U$ is contractible, there exists a smooth deformation retraction
    \begin{equation}
        r : U \times [0,1] \longrightarrow U,
    \end{equation}
    with $r(t,0) = t$ and $r(t,1) = t_0$.
    The composition
    \begin{equation}
        H(z,t,\tau) = \Psi_{t, r(t,\tau)}(z)
    \end{equation}
    defines a smooth isotopy in $M$ between the level sets $f_t^{-1}(c)$ and $f_{t_0}^{-1}(c)$ for all $t \in U$.

    This isotopy induces canonical isomorphisms
    \begin{equation}
        H^n(M, f_t^{\pm\infty}; \mathbb{C}) \cong H^n(M, f_{t_0}^{\pm\infty}; \mathbb{C}),
    \end{equation}
    varying smoothly with $t$.
    Hence the relative cohomology groups form a local system over $S$, and therefore determine a smooth vector bundle equipped with the flat Gauss--Manin connection.
    By Theorem~\ref{thm:main}, the rank of this bundle equals $\mu$.
\end{proof}

\begin{remark}\label{rem:picard-lefschetz}
    In classical Picard--Lefschetz theory, when $f$ is a Morse function the Lefschetz thimbles associated with its critical points form a basis of the relative homology
    \begin{equation}
        H_n\bigl(M, \{\operatorname{Re} f \ll 0\}\bigr),
    \end{equation}
    see for instance \cite{milnor1968singular}.
    This description relies on the Morse condition and therefore applies only to those deformation parameters for which the function $f_t$ has nondegenerate critical points.

    In contrast, the construction above does not require $f_t$ to be Morse.
    The relative homotopy type at infinity is identified directly with that of a regular fiber, so that the resulting relative cohomology bundle and its flat connection are defined over the whole deformation parameter space.
\end{remark}

% Retain the article's original bibliography selection from the shared library.
\nocite{arnold2012singularities,Batyrev1993,Broughton1988,CecottiVafa1991,CIR2014,Dubrovin1990ModernIII,dimca1992singularities,Dubrovin1993,Dubrovin1996,fan2011,FanFang2021,FanLanYang2020,FanJarvisRuan2013,GrauertRemmert2004,Hertling2003,Kulikov_1998,LeDungTrang1977,LI2022108165,Marsh2013TheBC,milnor1963morse,milnor1968singular,saito1989brieskorn,Saito1981Primitive,SaitoTakahashi2008,Wen2016,WenFan2015}
\printbibliography

\end{document}
